% page 510 du fac-simile (image p-509 du scan)
% bloc 157.5 x 237.7 mm ; marge gauche 8.0 mm
\pgc{-12.700mm}{0.000mm}{
\l{\cel{3}502.}
\l{}
\l{\sou{rolero}.\cel{2}Cilindro, ligna, metala, e c., quan onu promenas}
\l{\cel{3}sur agro por ruptar la glebi, por kunpresar tero ; por}
\l{\cel{3}planigar la gazon-agri ; por extensar pasto por la kukaji ;}
\l{\cel{3}por extensar la inko sur la komposturi di imprimilo, e c.}
\l{\cel{3}- EF.}
\l{}
\l{\sou{romano.}\cel{2}(I) Verko literaturala, fingura, en qua la autoro}
\l{\cel{3}esforcas interesar la lektero, sive\cel{2}per la singulareso}
\l{\cel{3}di la aventuri, sive per la deskripto di la mori, di la}
\l{\cel{3}pasioni, e c. - (II) Aventuri romanatra extraordinara.}
\l{\cel{3}- DEFIRS.}
\l{}
\l{\sou{romanco}. I. Mikra kansono di qua la karaktero esas sentimen-}
\l{\cel{3}tala. - II. Melodio qua havas tala karaktero. - DFIRS.}
\l{}
\l{\sou{romantika}. Qua apartenas a la sistemo literaturala qua refu-}
\l{\cel{3}zas la imito klasika di la Greki e di la Romani antiqua,}
\l{\cel{3}e volas liberigar su de la reguli kompozala qui konsequas}
\l{\cel{3}de ta imito.- DEFIRS.}
\l{}
\l{\sou{rombo}. (\sou{geom.)}\cel{2}Quadrilatero, kun lateri inter-egala e kun}
\l{\cel{3}anguli ne-orta.}
\l{}
\l{\sou{romboedro}. (\sou{geom.})\cel{2}Korpo solida, kristalajo, di qua la edri}
\l{\cel{3}esas romba. - DEF.}
\l{}
\l{\sou{ronda}. Di qua la formo esas cirkla. - DEFIS.}
\l{}
\l{\sou{rondelo}.\cel{2}I. (\sou{poezio}).Mikra poemo Franca, ek dek-e-tri versi}
\l{\cel{3}(ok kun ula rimo, e kin kun altra rimo) interseparita per}
\l{\cel{3}pauzo ye la verso kinesma, e ye la okesma - la vorti verso-}
\l{\cel{3}komencala iteresas pos la verso okesma e la dek-e-triesma,}
\l{\cel{3}sen esar parti de li. - II. (\sou{muziko}) Muzikajo di qua la}
\l{\cel{3}temato rivenas plurafoye. - DEFIRS.}
\l{}
\l{\sou{ronkar}. (\sou{netrans.}) Emisar, dormante, respiro nazala tre bruis-}
\l{\cel{3}oza. - FS.}
\l{}
\l{\sou{ronronar}. (\sou{netrans.)}\cel{2}(\sou{aludante kato}). Agar mikra grondado de}
\l{\cel{3}kontenteso. - F.}
\l{}
\l{\sou{roquar}. (\sou{netrans.}) (\sou{shakoludo)}. Pozar la turmo apud la rejo,}
\l{\cel{3}e la rejo trans la turmo kande nula piono stacas inter li,}
\l{\cel{3}o dextre, o sinistre, kande onu judikas lo kom oportuna. F.}
\l{}
\l{\sou{roso}. (\sou{meteorol}.) Strato de aquo-guteti qua aparas quale}
\l{\cel{3}sedimento, nokte, sur la surfaco di ula korpi qui expoz-}
\l{\cel{3}esas en aero ne-inkluzita - aquo-vaporo de atmosfero, kon-}
\l{\cel{3}densita da la varmesko quan produktas la radiado noktala.}
\l{\cel{3}- FRS.}
\l{}
\l{rosbifo.- (koquarto).Peco de bovokarno rostita - maxim-multa-}
\l{\cel{3}kaze pseudo-fileto. - DEFIRS.}
}
